\chapter*{Résumé}
\addcontentsline{toc}{chapter}{Résumé}

Ce projet informatique de Master 1 présente le développement d'un système d'estimation
automatique de la valeur des biens immobiliers résidentiels (maisons et appartements) au
Sénégal, à partir de techniques d'apprentissage automatique. L'objectif est double : d'une
part comparer plusieurs modèles de prédiction en optimisant systématiquement leurs
hyperparamètres, et d'autre part rendre ces résultats accessibles à travers une interface
graphique interactive.

En l'absence de jeu de données immobilier sénégalais existant, un corpus de 660 annonces de
vente a été collecté par extraction automatisée (\emph{web scraping}) depuis la plateforme
Expat-Dakar, couvrant l'ensemble du territoire national. Après nettoyage (suppression des prix
non exploitables et des biens hors périmètre résidentiel unitaire), un jeu de données de 603
annonces a été retenu, décrit par la surface, le nombre de chambres, le type de bien et le
quartier.

Quatre modèles ont été entraînés et comparés : une régression linéaire (référence), une
forêt aléatoire (\emph{Random Forest}), un modèle de \emph{gradient boosting} (XGBoost), et un
réseau de neurones multicouches (MLP) — les trois derniers optimisés par recherche
aléatoire d'hyperparamètres avec validation croisée. Le modèle Random Forest optimisé obtient
les meilleures performances sur le jeu de test (MAE de 73,3 millions FCFA, MAPE de 34,0~\%,
$R^2$ de 0,751), devançant XGBoost et le réseau de neurones — un résultat cohérent avec la
littérature sur des jeux de données tabulaires de taille modeste.

Une application interactive développée avec Streamlit permet enfin d'estimer le prix d'un bien
à partir de ses caractéristiques, d'explorer le jeu de données et de comparer visuellement les
performances des différents modèles.

\vspace{0.5cm}
\noindent\textbf{Mots-clés :} immobilier, apprentissage automatique, optimisation
d'hyperparamètres, Random Forest, XGBoost, réseau de neurones, Sénégal, Streamlit.
